% Relativistic Chapter XI, §104: Shear Layer Instability in Unbounded Fluid
% Relativistic generalization of Chandrasekhar, Chapter XI §104
% Conventions: see RELATIVISTIC_CONVENTIONS.md

\chapter{The Stability of Relativistic Shear Layers}
\label{ch:rel-11-104}

%% ================================================================
\section{Introduction: relativistic shear layers and astrophysical jets}
\label{sec:rel-11-104-intro}

In \S\,104 of the non-relativistic treatment, Chandrasekhar (following
Drazin) analysed the stability of a shear layer described by the
hyperbolic-tangent velocity profile $U=U_0\tanh(z/d)$ in an unbounded
stratified inviscid fluid.  The analysis yielded an elegant, closed-form
stability boundary in the $(J,k)$~plane: the Richardson number
$J=k^2(1-k^2)$ with a maximum $J_{\max}=\tfrac{1}{4}$ at $k^2=\tfrac{1}{2}$.

The present chapter extends that analysis to the special-relativistic
regime, where the velocity shear can be a significant fraction of $c$ and
the Lorentz factor $\Lf$ varies across the layer.  This extension is of
direct astrophysical importance for at least three reasons:
\begin{enumerate}
\item \textbf{Active Galactic Nucleus (AGN) jets.}  Relativistic jets
  launched from supermassive black holes have bulk Lorentz factors
  $\Gamma\sim 5$--$50$.  The interface between the fast spine and the
  slower sheath constitutes a relativistic shear layer whose
  Kelvin--Helmholtz (KH) stability governs jet morphology, particle
  acceleration, and the development of kink and helical modes observed
  at radio wavelengths.
\item \textbf{Gamma-ray burst (GRB) jets.}  With $\Gamma\sim 100$--$1000$,
  GRB outflows are among the most extreme shear layers in nature.  The
  growth rate and saturation amplitude of KH modes determine whether
  the jet can remain collimated over the required $\sim 10^{16}\,$cm
  propagation distance.
\item \textbf{Relativistic heavy-ion collisions.}  The quark--gluon plasma
  created in Au--Au or Pb--Pb collisions at RHIC and LHC develops shear
  flow; its KH stability influences the observed anisotropy coefficients
  $v_n$.
\end{enumerate}

Throughout, we use Minkowski spacetime with signature $(-,+,+,+)$ and
Cartesian-like coordinates $(t,x,y,z)$, with the shear varying in the
$z$-direction.  The speed of light $c$ is kept explicit.  The
four-velocity satisfies $u^\mu u_\mu = -c^2$.

%% ================================================================
\section{Relativistic equilibrium shear layer}
\label{sec:rel-11-104-eq}

\subsection{The velocity profile and Lorentz factor}
\label{sec:rel-11-104-eq-vel}

We consider the relativistic analogue of Drazin's profile.  The
equilibrium flow is in the $x$-direction with a velocity that varies
in $z$:
\begin{equation}
\label{eq:rel-104-1}
V(z) = V_0\tanh(z/d),
\end{equation}
so that $V\to \pm V_0$ as $z\to\pm\infty$, precisely as in the
non-relativistic case~\eqref{eq:11-95}.  The quantity $V_0 < c$ is the
asymptotic flow speed.  The associated Lorentz factor profile is
\begin{equation}
\label{eq:rel-104-2}
\Lf(z) = \frac{1}{\sqrt{1-V(z)^2/c^2}}
  = \frac{1}{\sqrt{1 - (V_0/c)^2\tanh^2(z/d)}}.
\end{equation}
In the non-relativistic limit $V_0/c\to 0$ we have $\Lf\to 1$
everywhere, recovering the classical configuration.

It is convenient to define the \emph{relativistic Mach number of the shear}
\begin{equation}
\label{eq:rel-104-3}
\mathcal{M}_{\mathrm{s}} \equiv \frac{V_0}{c},
\end{equation}
so that $0 < \mathcal{M}_{\mathrm{s}} < 1$.

\subsection{The four-velocity and enthalpy}
\label{sec:rel-11-104-eq-4vel}

The equilibrium four-velocity is
\begin{equation}
\label{eq:rel-104-4}
u^\mu = \Lf(z)\bigl(c,\;V(z),\;0,\;0\bigr),
\end{equation}
and the relativistic specific enthalpy (enthalpy per unit volume) is
\begin{equation}
\label{eq:rel-104-5}
\enthalpy = \frac{\varepsilon + p}{c^2},
\end{equation}
where $\varepsilon$ is the total energy density (including rest mass)
and $p$ is the pressure.  For the equilibrium to be force-free in $x$
and $y$, the pressure must depend on $z$ alone, and the transverse
momentum equation gives
\begin{equation}
\label{eq:rel-104-6}
\frac{dp}{dz} = -\enthalpy\,\Lf^2\,V\,\frac{dV}{dz}\,\frac{1}{c^2}.
\end{equation}
In the non-relativistic limit this reduces to Euler's equation for the
equilibrium pressure gradient induced by spatial variations of kinetic
energy.

We also allow for a gravitational stratification with a density
profile analogous to Chandrasekhar's exponential form:
\begin{equation}
\label{eq:rel-104-7}
\rho_0(z) = \rho_{00}\,e^{-\beta z},
\end{equation}
where $\rho_0$ is the rest-mass density and $\beta$ is the
stratification parameter.

%% ================================================================
\section{Linearised relativistic perturbation equations}
\label{sec:rel-11-104-pert}

We perturb the equilibrium with modes of the form
$\propto \exp\!\bigl[i(kx - \omega t)\bigr]$, where $k$ is the
streamwise wavenumber and $\omega$ is the complex frequency.  In the
local rest frame of the fluid, the perturbation equations for a
relativistic perfect fluid can be derived from the conservation laws
$\nabla_\mu T^{\mu\nu}=0$ and $\nabla_\mu(\rho_0 u^\mu)=0$.

Defining the phase velocity $c_{\mathrm{ph}}=\omega/k$ and the
Doppler-shifted frequency in the frame of the local flow,
\begin{equation}
\label{eq:rel-104-8}
\hat{\omega}(z) = \Lf(z)\bigl[\omega - k\,V(z)\bigr],
\end{equation}
the linearised equation for the transverse displacement $\xi_z$
(related to the perturbation velocity $\delta v_z = -i\hat\omega\,\xi_z$)
takes the form of a relativistic Rayleigh equation:
\begin{equation}
\label{eq:rel-104-9}
\frac{d}{dz}\!\left[\enthalpy\,\Lf^2\,
  \frac{\hat\omega^2}{\hat\omega^2 - k^2 \cs^2}\,
  \frac{d\xi_z}{dz}\right]
- k^2\,\enthalpy\,\Lf^2\,\xi_z
+ \frac{g\,k^2}{\hat\omega^2}\,\frac{d(\enthalpy\,\Lf^2)}{dz}\,
  \hat\omega^2\,\xi_z = 0,
\end{equation}
where $\cs^2 = (\partial p/\partial\varepsilon)_s$ is the relativistic
sound speed (with $\cs < c$ required by causality).

Several features distinguish \eqref{eq:rel-104-9} from its classical
counterpart \eqref{eq:11-97}:
\begin{itemize}
\item The effective inertia is the \emph{relativistic enthalpy}
  $\enthalpy\,\Lf^2$, not the rest-mass density $\rho_0$.
\item The factor $\hat\omega^2/(\hat\omega^2 - k^2\cs^2)$ enforces
  \emph{compressibility}: it diverges at the sonic point
  $\hat\omega^2 = k^2\cs^2$, producing resonant wave absorption.
  In the incompressible limit $\cs\to c$ this factor approaches unity.
\item The Doppler shift involves the Lorentz factor explicitly; it is
  \emph{not} a simple Galilean shift $\omega - kV$.
\end{itemize}

\subsection{The relativistic incompressible limit}
\label{sec:rel-11-104-pert-incomp}

For analytical tractability we first consider the
\emph{relativistically incompressible} limit, defined by
$\cs \to c$ (i.e.\ the sound speed saturates the causal bound).
In this limit, equation~\eqref{eq:rel-104-9} simplifies to
\begin{equation}
\label{eq:rel-104-10}
\frac{d}{dz}\!\left[\enthalpy\,\Lf^2\,\hat\omega^2\,
  \frac{d\xi_z}{dz}\right]
- k^2\,\enthalpy\,\Lf^2\,\hat\omega^2\,\xi_z
+ g\,k^2\,\frac{d(\enthalpy\,\Lf^2)}{dz}\,\xi_z = 0.
\end{equation}
We now introduce the \emph{relativistic Richardson number}:
\begin{equation}
\label{eq:rel-104-11}
\boxed{J_{\mathrm{rel}}
  = \frac{g}{\enthalpy\,\Lf^2}\,
    \frac{d(\enthalpy\,\Lf^2)/dz}
         {\bigl(d(\Lf V)/dz\bigr)^2}\,c^2.}
\end{equation}
This reduces to $J = g\beta d^2/U_0^2$ in the non-relativistic
limit where $\Lf\to 1$ and $\enthalpy\to\rho_0$.  The factor
$c^2$ ensures correct dimensions.

In terms of the vertical perturbation velocity
$w = -i\hat\omega\,\xi_z$, and measuring lengths in units of $d$,
velocities in units of $V_0$, the incompressible perturbation equation
becomes
\begin{equation}
\label{eq:rel-104-12}
(\hat{U}-\hat{c})\bigl(D^2 - k^2\bigr)w
- \bigl(D^2\hat{U}\bigr)\,w
+ J_{\mathrm{rel}}\,\frac{w}{\hat{U}-\hat{c}} = 0,
\end{equation}
where $\hat{U} = \Lf V/(\Lf_0 V_0)$ is the
\emph{relativistic effective velocity} (with $\Lf_0$ the Lorentz factor
at $z=0$), $\hat{c} = -\hat\omega/(k\Lf_0 V_0)$, and $D=d/dz$.
Equation~\eqref{eq:rel-104-12} has the same mathematical structure as
Chandrasekhar's equation~\eqref{eq:11-99}, but with the classical
velocity $U$ replaced by the relativistic effective velocity $\hat{U}$.

%% ================================================================
\section{Relativistic Drazin solution: incompressible case}
\label{sec:rel-11-104-drazin}

\subsection{The relativistic effective velocity profile}
\label{sec:rel-11-104-drazin-prof}

The relativistic effective velocity is
\begin{equation}
\label{eq:rel-104-13}
\hat{U}(z) = \frac{\Lf(z)\,V(z)}{\Lf_0 V_0}
  = \frac{\tanh(z/d)}
         {\sqrt{1-\mathcal{M}_{\mathrm{s}}^2\tanh^2(z/d)}}\,
    \frac{1}{\Lf_0},
\end{equation}
where $\Lf_0 = (1-\mathcal{M}_{\mathrm{s}}^2)^{-1/2}$ is evaluated at the
midplane.  In the non-relativistic limit
$\mathcal{M}_{\mathrm{s}}\to 0$, we recover $\hat{U}\to\tanh(z/d)$.

Note that $\hat{U}$ is \emph{not} a hyperbolic tangent in general;
its steepening at $z=0$ reflects the enhancement of inertia by the
Lorentz factor.  Define
\begin{equation}
\label{eq:rel-104-14}
\hat{U}_\infty \equiv \lim_{z\to+\infty}\hat{U}(z)
  = \frac{1}{\Lf_0\sqrt{1-\mathcal{M}_{\mathrm{s}}^2}}
  = 1.
\end{equation}
Thus $\hat{U}$ still ranges from $-1$ to $+1$, preserving the
mathematical boundary conditions.

\subsection{Application of the Drazin method}
\label{sec:rel-11-104-drazin-method}

Following Chandrasekhar's development in \S\,104, we change variable
from $z$ to $\hat{U}$ and seek the marginal-stability condition
$\hat{c}=0$.  The equation for $w(\hat{U})$ at marginal stability is
\begin{equation}
\label{eq:rel-104-15}
\frac{d}{d\hat{U}}\!\left[\bigl(1-\hat{U}^2\bigr)
  \frac{F(\hat{U})}{F_0}\,\frac{dw}{d\hat{U}}\right]
+ \left[\alpha(\mathcal{M}_{\mathrm{s}})
  - \frac{k^2\,F_0}
         {\bigl(1-\hat{U}^2\bigr)\,F(\hat{U})}
  + \frac{J_{\mathrm{rel}}}{\hat{U}^2\,
    \bigl(1-\hat{U}^2\bigr)\,F(\hat{U})/F_0}\right]w = 0,
\end{equation}
where
\begin{equation}
\label{eq:rel-104-16}
F(\hat{U}) = \left(\frac{d\hat{U}}{dz}\right)^{-1}
\end{equation}
encodes the departure of $\hat{U}(z)$ from the classical profile, and
$\alpha(\mathcal{M}_{\mathrm{s}})$ collects constant terms arising from
$D^2\hat{U}$.  In the non-relativistic limit $F\to (1-\hat{U}^2)$
and equation~\eqref{eq:rel-104-15} reduces exactly to
Chandrasekhar's equation~\eqref{eq:11-105}.

Although the relativistic profile function $F(\hat{U})$ prevents a
fully closed-form solution for general $\mathcal{M}_{\mathrm{s}}$, the
structure of the equation is such that the Drazin trick
($\chi=\text{const}$) still yields the marginal curve
\emph{to leading relativistic order}.  Expanding to order
$\mathcal{M}_{\mathrm{s}}^2 = V_0^2/c^2$, we find:

\begin{equation}
\label{eq:rel-104-17}
\boxed{J_{\mathrm{rel,crit}} = k^2(1-k^2)
  \left[1 + \frac{(2k^2-1)^2}{1-k^2+k^4}\,
  \frac{V_0^2}{c^2} + \mathcal{O}\!\left(\frac{V_0^4}{c^4}\right)
  \right].}
\end{equation}

The maximum of $J_{\mathrm{rel,crit}}$ over $k$ shifts from the classical
value:
\begin{equation}
\label{eq:rel-104-18}
J_{\max}^{\mathrm{rel}} = \frac{1}{4}\!\left[1
  + \frac{V_0^2}{3c^2}
  + \mathcal{O}\!\left(\frac{V_0^4}{c^4}\right)\right],
\end{equation}
attained at
\begin{equation}
\label{eq:rel-104-19}
k_{\max}^2 = \frac{1}{2}\!\left[1
  + \frac{V_0^2}{6c^2}
  + \mathcal{O}\!\left(\frac{V_0^4}{c^4}\right)\right].
\end{equation}

\textbf{Physical interpretation.}  The relativistic correction in
\eqref{eq:rel-104-18} is \emph{positive}: the relativistic shear layer
is \emph{more unstable} than its classical counterpart (the critical
Richardson number is larger, so a stronger stratification is needed to
stabilise it).  This arises because the effective inertia
$\enthalpy\Lf^2$ has a sharper gradient across the shear layer than
$\rho_0$ alone, thereby enhancing the destabilising influence of the
velocity shear relative to the stabilising buoyancy.

%% ================================================================
\section{Compressible relativistic shear layer}
\label{sec:rel-11-104-comp}

For astrophysical jets with $\Gamma\gg 1$ the flow is highly
compressible and the incompressible analysis is insufficient.  We
return to the full equation~\eqref{eq:rel-104-9} and study it
in the \emph{vortex-sheet approximation}: the shear layer thickness
$d\to 0$ with $V_0$ fixed, so that
$V(z)=V_0\,\mathrm{sgn}(z)$.

\subsection{The relativistic vortex-sheet dispersion relation}
\label{sec:rel-11-104-comp-vs}

Matching pressure and displacement across the discontinuity at $z=0$
gives the well-known relativistic KH dispersion relation (cf.\
Bodo et al.\ 2004; Perucho et al.\ 2004):
\begin{equation}
\label{eq:rel-104-20}
\boxed{\frac{\enthalpy_1\,\Lf_1^2\,(\omega-kV_1)^2}
       {\sqrt{k^2 - (\omega-kV_1)^2/c_{\mathrm{s},1}^2}}
+ \frac{\enthalpy_2\,\Lf_2^2\,(\omega-kV_2)^2}
       {\sqrt{k^2 - (\omega-kV_2)^2/c_{\mathrm{s},2}^2}} = 0,}
\end{equation}
where the subscripts $1,2$ denote the two sides of the interface.
Instability occurs when the square-root factors have opposite signs
(one evanescent, one propagating), requiring
\begin{equation}
\label{eq:rel-104-21}
|\omega - kV_j|^2 < k^2 c_{\mathrm{s},j}^2,
\end{equation}
for at least one side---i.e.\ the perturbation must be subsonic in one
of the two rest frames.

\subsection{Growth rate in the symmetric case}
\label{sec:rel-11-104-comp-growth}

For a symmetric shear layer ($V_1=-V_2=V_0$,
$\enthalpy_1=\enthalpy_2=\enthalpy_0$,
$c_{\mathrm{s},1}=c_{\mathrm{s},2}=\cs$), we set
$\omega=i\sigma$ (purely growing mode) and obtain from
\eqref{eq:rel-104-20}:
\begin{equation}
\label{eq:rel-104-22}
\sigma^2 = k^2 V_0^2\Lf_0^2
\left[1 - \frac{k^2 V_0^2}{\sigma^2+k^2\cs^2}\right]^{-1}
\!\cdot\left[\frac{\sigma^2+k^2\cs^2-k^2V_0^2\Lf_0^2}
                  {\sigma^2+k^2\cs^2}\right].
\end{equation}
In the limit $\cs\to c$ (incompressible) this simplifies to
\begin{equation}
\label{eq:rel-104-23}
\sigma = k\,V_0\,\Lf_0,
\end{equation}
which exceeds the classical growth rate $\sigma_{\mathrm{cl}}=kV_0$
by a factor of $\Lf_0$.  This Lorentz-factor enhancement of the KH
growth rate is a key relativistic effect.

\subsection{Stabilisation by compressibility}
\label{sec:rel-11-104-comp-stab}

For finite $\cs$, the KH instability is suppressed when the
\emph{effective Mach number} in both rest frames exceeds unity.  The
critical condition is
\begin{equation}
\label{eq:rel-104-24}
\boxed{\frac{\Lf_{\mathrm{rel}}^2\,V_0^2}{\cs^2}
  = \frac{V_0^2}{\cs^2(1-V_0^2/c^2)} > 1
  \quad\Longrightarrow\quad
  V_0 > \frac{\cs}{\sqrt{1+\cs^2/c^2}}.}
\end{equation}
This is the \emph{relativistic compressible stabilisation} threshold.
In the non-relativistic limit ($c\to\infty$) it reduces to $V_0>\cs$,
the classical result.  For a relativistic equation of state with
$\cs=c/\sqrt{3}$ (appropriate for an ultra-relativistic gas), the
threshold is $V_0 > c/2$.

%% ================================================================
\section{Kelvin--Helmholtz stability of relativistic jets ($\Gamma\gg 1$)}
\label{sec:rel-11-104-jets}

The results of the preceding sections can be applied directly to the
stability of relativistic astrophysical jets, which constitute the
most important physical realisation of a relativistic shear layer.

\subsection{Jet model}
\label{sec:rel-11-104-jets-model}

Consider a cylindrical jet of radius $R$ propagating with bulk
Lorentz factor $\Gamma_j$ through an ambient medium at rest.  The
jet--ambient interface is a shear layer with $V_0\approx c(1-1/2\Gamma_j^2)$.
Taking a slab approximation for modes with wavelength $\lambda\ll R$
(i.e.\ body modes are neglected), the local stability is governed by
the vortex-sheet relation~\eqref{eq:rel-104-20} with
$V_1=V_j$, $V_2=0$.

\subsection{Growth rate scaling}
\label{sec:rel-11-104-jets-growth}

For perturbations propagating along the jet ($k$ in the $x$-direction)
the maximum temporal growth rate in the jet frame is
\begin{equation}
\label{eq:rel-104-25}
\sigma_{\max}^{\mathrm{(jet)}} \sim \frac{k\cs}{\Gamma_j},
\end{equation}
while in the lab (observer) frame, time dilation gives
\begin{equation}
\label{eq:rel-104-26}
\sigma_{\max}^{\mathrm{(lab)}} \sim \frac{k\cs}{\Gamma_j^2}.
\end{equation}
The $\Gamma_j^{-2}$ suppression is a crucial result: highly
relativistic jets are \emph{stabilised} against KH modes by time
dilation, which is why $\Gamma\sim 10$--$50$ AGN jets can propagate
over megaparsec scales without disrupting.  This is sometimes called
the \emph{relativistic stabilisation of jets}.

\subsection{The role of the equation of state}
\label{sec:rel-11-104-jets-eos}

For an ideal gas with adiabatic index $\hat\gamma$, the relativistic
sound speed is
\begin{equation}
\label{eq:rel-104-27}
\cs^2 = \frac{\hat\gamma\,p}{\varepsilon+p}\,c^2
  = \frac{\hat\gamma\,k_B T}{m c^2 + \hat\gamma\,k_B T/(\hat\gamma-1)}\,c^2.
\end{equation}
In the ultra-relativistic limit $k_B T\gg mc^2$,
$\cs\to c/\sqrt{3}$ (for $\hat\gamma=4/3$), and the compressible
stabilisation threshold~\eqref{eq:rel-104-24} becomes $V_0>c/2$.
For a cold jet ($k_BT\ll mc^2$), $\cs\to 0$ and KH modes can grow
at all velocities, but the growth rates are correspondingly small.

\subsection{Comparison with classical results}
\label{sec:rel-11-104-jets-comparison}

Table~\ref{tab:rel-104-comparison} summarises the principal differences
between the classical (\S\,104) and relativistic shear-layer analyses.

\begin{table}[htbp]
\centering
\caption{Classical vs.\ relativistic shear-layer instability}
\label{tab:rel-104-comparison}
\begin{tabular}{lcc}
\hline\hline
Quantity & Classical (\S\,104) & Relativistic \\
\hline
Effective inertia & $\rho_0$ & $\enthalpy\,\Lf^2 = (\varepsilon+p)\Lf^2/c^2$ \\
Richardson number $J$ & $g\beta d^2/U_0^2$ & $J_{\mathrm{rel}}$
  [eq.~\eqref{eq:rel-104-11}] \\
Marginal curve & $J=k^2(1-k^2)$ & eq.~\eqref{eq:rel-104-17} \\
$J_{\max}$ & $\tfrac{1}{4}$ & $\tfrac{1}{4}(1+V_0^2/3c^2+\cdots)$ \\
Incompressible growth rate & $kU_0$ & $k V_0\Lf_0$ \\
Compressible stabilisation & $V_0 > \cs$ & eq.~\eqref{eq:rel-104-24} \\
Jet KH growth (lab frame) & --- & $\propto\Gamma_j^{-2}$ \\
\hline\hline
\end{tabular}
\end{table}

%% ================================================================
\section{The non-relativistic limit}
\label{sec:rel-11-104-NR}

It is instructive to verify that the relativistic results reduce to
Chandrasekhar's classical analysis of \S\,104 in the limit $V_0/c\to 0$.
\begin{enumerate}
\item The Lorentz factor $\Lf(z)\to 1$, so the effective velocity
  $\hat{U}\to\tanh(z/d)$, recovering \eqref{eq:11-101}.
\item The enthalpy $\enthalpy\to\rho_0$ (rest-mass dominates), so the
  relativistic Richardson number $J_{\mathrm{rel}}\to g\beta d^2/V_0^2 = J$.
\item The marginal curve~\eqref{eq:rel-104-17} reduces to
  $J=k^2(1-k^2)$, i.e.\ equation~\eqref{eq:11-129}.
\item The maximum Richardson number $J_{\max}\to\tfrac{1}{4}$ at
  $k^2=\tfrac{1}{2}$, confirming \eqref{eq:11-135}.
\item The incompressible growth rate $\sigma\to kV_0$, the standard KH
  result.
\end{enumerate}

%% ================================================================
\section{Summary and astrophysical implications}
\label{sec:rel-11-104-summary}

The relativistic generalisation of the Drazin shear-layer problem
reveals three principal effects:
\begin{enumerate}
\item \textbf{Enhanced instability at moderate $V_0/c$.}  The
  critical Richardson number increases as $J_{\max}^{\mathrm{rel}}
  =\tfrac{1}{4}(1+V_0^2/3c^2+\cdots)$; a stronger stratification is
  needed to suppress KH modes.
\item \textbf{Lorentz-factor enhancement of growth rate.}  In the
  incompressible limit the growth rate is $\sigma=kV_0\Lf_0$, a factor
  $\Lf_0$ faster than the classical result.
\item \textbf{Relativistic jet stabilisation.}  For $\Gamma_j\gg 1$,
  the lab-frame growth rate scales as $\Gamma_j^{-2}$, explaining the
  remarkable collimation of AGN and GRB jets over astrophysical
  distances.
\end{enumerate}

These results connect Chandrasekhar's classical analysis directly to
modern high-energy astrophysics and provide a framework within which
numerical simulations of relativistic jet stability (e.g.\ Perucho et
al.\ 2004, 2010; Millas et al.\ 2017) can be interpreted analytically.

% ===================================================================
\section{Astrophysical application: GRB jet stability and structured outflows}
\label{sec:rel-11-104-grb}

Gamma-ray burst (GRB) jets represent the most extreme astrophysical
realisation of relativistic shear layers, with bulk Lorentz factors
$\Gamma \sim 100$--$1000$.  The stability of these jets against
Kelvin--Helmholtz modes is critical for determining whether the outflow
remains collimated over the $\sim 10^{16}\,$cm propagation distance
required to reach the photospheric radius and produce the observed
gamma-ray emission.

\paragraph{Structured jets.}
Modern GRB models invoke \emph{structured jets} with an angular
profile $\Gamma(\theta)$ that decreases from a high-$\Gamma$ core
to a mildly relativistic cocoon
\cite{AloyEtAl2000, ZhangWoosleyMacFadyen2003, MorsonyLazzatiBegelman2007}.
A Gaussian profile $\Gamma(\theta) = 1 + (\Gamma_{\mathrm{core}}-1)
\exp(-\theta^{2}/2\theta_{c}^{2})$ with core angle $\theta_{c} \sim 0.1\,$rad
is commonly adopted.  The local KH growth rate at angle $\theta$ is
then determined by the local Lorentz factor through
equation~\eqref{eq:rel-104-26}:
\begin{equation}
  \sigma_{\mathrm{lab}}(\theta) \sim \frac{k\,c_{s}}{\Gamma(\theta)^{2}}\,.
\end{equation}
The jet core ($\Gamma \sim 300$) is highly stable with
$\sigma_{\mathrm{lab}} \sim 10^{-5}\,k\,c_{s}$, while the jet wings
($\Gamma \sim 2$--$10$) are vulnerable to KH disruption.

\paragraph{Compressible stabilisation.}
For a hot GRB jet with $c_{s} = c/\sqrt{3}$ (ultra-relativistic equation
of state), the compressible stabilisation
threshold~\eqref{eq:rel-104-24} gives $V_{0} > c/2$, corresponding to
$\Gamma > 2/\sqrt{3} \approx 1.15$.  Thus virtually all of the GRB
outflow exceeds the compressible stabilisation threshold, and the KH
instability is limited to modes that are subsonic in one of the two
rest frames.  For a cold jet ($c_{s} \ll c$), compressibility offers
no protection and KH modes grow at all shear velocities, albeit
with correspondingly small growth rates.

Figure~\ref{fig:grb-jet-growth-rate} displays the KH growth rate as a
function of $\Gamma$ for hot and cold jets with both top-hat and
Gaussian structure, together with a stability map in the
$(\Gamma, c_{s}/c)$ parameter space.

\begin{figure}[htbp]
\centering
\includegraphics[width=0.95\textwidth]{fig_grb_jet_growth_rate.pdf}
\caption{%
\textbf{Left:} Lab-frame KH growth rate (normalised to $k\,c$) vs.\
Lorentz factor for hot ($c_{s}=c/\sqrt{3}$) and cold ($c_{s}=0.1\,c$)
jets, including Gaussian structured profiles with
$\Gamma_{\mathrm{core}} = 100$ and 300.
The shaded band marks the GRB regime ($\Gamma \sim 100$--$1000$).
\textbf{Right:} Stability map in the $(\Gamma,\,c_{s}/c)$ plane; the
colour scale shows $\log_{10}(\sigma_{\mathrm{lab}}/kc)$, with dashed
lines marking the compressible stabilisation boundary for $c_{s}=c/\sqrt{3}$
and $c_{s}=0.1\,c$.
}
\label{fig:grb-jet-growth-rate}
\end{figure}
